\documentclass[11pt,a4paper]{article}
% preamble.tex — shared setup for human-verification documents
% Compile each document with:  tectonic <File>.tex

\usepackage[T1]{fontenc}
\usepackage[utf8]{inputenc}
\usepackage{lmodern}
\usepackage{amssymb,amsmath,amsthm}
\usepackage{mathtools}
\usepackage[dvipsnames]{xcolor}
\usepackage{listings}
\usepackage{hyperref}
\usepackage[margin=2.5cm]{geometry}
\usepackage{parskip}

% ── Colours ──────────────────────────────────────────────────────────
\definecolor{leanBg}{HTML}{F7F7F0}
\definecolor{haskellBg}{HTML}{F0F4F7}
\definecolor{keyword}{HTML}{7B3294}
\definecolor{comment}{HTML}{5A8C32}
\definecolor{noteFg}{HTML}{B35900}

% ── Lean 4 listings language ─────────────────────────────────────────
% NOTE: Lean listing blocks use ASCII approximations of unicode symbols
% (e.g. -> for arrow, o. for compose).  The literate key substitutes
% these with proper math-mode renderings in the PDF.
\lstdefinelanguage{Lean4}{
  morekeywords={def,theorem,structure,where,namespace,end,import,fun,
                let,have,show,by,exact,intro,cases,rfl,subst,obtain,
                funext,simp,rw,apply,variable,Type,Prop,instance,
                section,open,class,extends,deriving,sorry,noncomputable},
  sensitive=true,
  morecomment=[l]{--},
  morecomment=[n]{/-}{-/},
  morestring=[b]",
}

% ── Lean listing style ───────────────────────────────────────────────
\lstdefinestyle{lean}{
  language=Lean4,
  backgroundcolor=\color{leanBg},
  basicstyle=\small\ttfamily,
  keywordstyle=\color{keyword}\bfseries,
  commentstyle=\color{comment}\itshape,
  breaklines=true,
  frame=single,
  framerule=0.4pt,
  xleftmargin=4pt,
  xrightmargin=4pt,
  aboveskip=6pt,
  belowskip=4pt,
  literate=
    {->}{{$\to$}}2
    {=>}{{$\Rightarrow$}}2
    {<|}{{$\langle$}}1
    {|>}{{$\rangle$}}1
    {==>}{{$\Longrightarrow$}}3
    { o. }{{$\;\circ\;$}}3
    {forall}{{$\forall$}}1
    {:=}{{$\,:=\,$}}2
    {!=}{{$\neq$}}2
    {<=}{{$\leq$}}2
    {>=}{{$\geq$}}2,
}

% ── Haskell listing style ───────────────────────────────────────────
\lstdefinestyle{haskell}{
  language=Haskell,
  backgroundcolor=\color{haskellBg},
  basicstyle=\small\ttfamily,
  keywordstyle=\color{keyword}\bfseries,
  commentstyle=\color{comment}\itshape,
  breaklines=true,
  frame=single,
  framerule=0.4pt,
  xleftmargin=4pt,
  xrightmargin=4pt,
  aboveskip=6pt,
  belowskip=4pt,
  literate=
    {->}{{$\to$}}2
    {=>}{{$\Rightarrow$}}2,
}

% ── Theorem environments ────────────────────────────────────────────
\theoremstyle{definition}
\newtheorem{defn}{Definition}[section]
\newtheorem{thm}[defn]{Theorem}
\newtheorem{prop}[defn]{Proposition}

\theoremstyle{remark}
\newtheorem*{fidelitynote}{Fidelity note}

% ── Convenience macros ───────────────────────────────────────────────
\newcommand{\lean}[1]{\texttt{#1}}      % inline Lean identifier
\newcommand{\hask}[1]{\texttt{#1}}      % inline Haskell identifier
\newcommand{\Ext}{\mathrm{Ext}}
\newcommand{\Pos}{\mathrm{Pos}}
\newcommand{\Shape}{\mathrm{Shape}}


\title{Human Verification: \texttt{Containers.Rungs}\\[4pt]
  \large DAG and Graph Rungs of the Trust Ladder}
\author{Machine-generated from \texttt{lean/Rungs.lean} for review by Neil Ghani}
\date{July 2026}

\begin{document}
\maketitle

\begin{abstract}
This document presents every definition and theorem \emph{statement} from the
Lean~4 module \lean{Containers.Rungs} in three parallel forms: the original
Lean, a readable Haskell rendering, and precise mathematical English.  Proofs
are omitted --- they are machine-checked by Lean.  The purpose is to let a
human verifier confirm that the formalisation captures the intended
mathematics.

The file verifies two negative-result rungs of the directed-container trust
ladder (\S7 and \S8 of the paper): the DAG rung shows that the Sawasdee
Cellars wine graph extended with a Singapore transshipment is \emph{not thin}
(two distinct routes between the same nodes exist), and the Graph rung shows
that a 2-node scrap loop produces an \emph{infinite} position set.
\end{abstract}

\tableofcontents
\bigskip

%===================================================================
\section{Overview}
%===================================================================

A directed container whose position sets are all singletons is called
\emph{thin}: there is at most one route between any two nodes.  The ladder in
the paper has four rungs:

\begin{enumerate}
  \item \textbf{Tree} --- positions are singletons (trivially thin).
  \item \textbf{Polytree / DAG without multiple routes} --- still thin.
  \item \textbf{DAG with multiple routes} --- not thin ($\lnot$ Subsingleton).
  \item \textbf{General graph with cycles} --- position sets are infinite.
\end{enumerate}

This file addresses rungs~3 and~4 with concrete witnesses.

%===================================================================
\section{DAG rung: Sawasdee + Singapore transshipment (\S7)}
%===================================================================

% --- Facility -------------------------------------------------------
\subsection{Inductive type: \lean{SawasdeeDAG.Facility}}

\begin{defn}[Extended facility set]
The set of \textbf{facilities} in the extended Sawasdee Cellars network consists
of the original eleven nodes plus Singapore: twelve facilities in total.
\end{defn}

\paragraph{Lean.}
\begin{lstlisting}[style=lean]
inductive Facility
  | vineyard1 | vineyard2 | vineyard3
  | sawasdee | melbournePort | singapore | thaiPort
  | distributor1 | distributor2
  | retail1 | retail2 | retail3
  deriving DecidableEq, Repr
\end{lstlisting}

\paragraph{Haskell.}
\begin{lstlisting}[style=haskell]
data Facility
  = Vineyard1 | Vineyard2 | Vineyard3
  | Sawasdee | MelbournePort | Singapore | ThaiPort
  | Distributor1 | Distributor2
  | Retail1 | Retail2 | Retail3
  deriving (Eq, Show)
\end{lstlisting}


% --- Leg -----------------------------------------------------------
\subsection{Inductive type: \lean{SawasdeeDAG.Leg}}

\begin{defn}[Extended leg set]
The \textbf{legs} of the extended network consist of the original ten shipping
legs plus two Singapore transshipment legs: \lean{sing1} (Melbourne Port
$\to$ Singapore) and \lean{sing2} (Singapore $\to$ Thai Port).
\end{defn}

\paragraph{Lean.}
\begin{lstlisting}[style=lean]
inductive Leg
  | v1s | v2s | v3s
  | smp
  | sea        -- direct Melbourne Port -> Thai Port
  | sing1      -- Melbourne Port -> Singapore
  | sing2      -- Singapore -> Thai Port
  | td1 | td2
  | d1r1 | d1r2
  | d2r3
  deriving DecidableEq, Repr
\end{lstlisting}

\paragraph{Haskell.}
\begin{lstlisting}[style=haskell]
data Leg
  = V1s | V2s | V3s | Smp
  | Sea | Sing1 | Sing2
  | Td1 | Td2 | D1r1 | D1r2 | D2r3
  deriving (Eq, Show)
\end{lstlisting}


% --- G' -------------------------------------------------------------
\subsection{Definition: \lean{SawasdeeDAG.G'}}

\begin{defn}[Extended Sawasdee graph]
$G'$ is the digraph with vertices \lean{Facility}, edges \lean{Leg}, and
source/target maps \lean{leg\_src} and \lean{leg\_tgt} defined by pattern
matching.  It extends the original Sawasdee Cellars graph with the Singapore
transshipment leg.
\end{defn}

\paragraph{Lean.}
\begin{lstlisting}[style=lean]
@[reducible] def G' : Digraph where
  V := Facility
  E := Leg
  src := leg_src
  tgt := leg_tgt
\end{lstlisting}

\paragraph{Haskell.}
\begin{lstlisting}[style=haskell]
data Digraph v e = Digraph { src :: e -> v, tgt :: e -> v }

g' :: Digraph Facility Leg
g' = Digraph legSrc legTgt
\end{lstlisting}


% --- directPath / transshipPath ------------------------------------
\subsection{Definitions: \lean{directPath} and \lean{transshipPath}}

\begin{defn}[Witness routes]
Two distinct directed paths from \lean{melbournePort} to \lean{thaiPort} in
$G'$:
\begin{itemize}
  \item $\text{directPath}$: the single leg $\mathrm{sea}$, of length~1.
  \item $\text{transshipPath}$: the two-leg route $\mathrm{sing1}$ followed by
    $\mathrm{sing2}$, of length~2.
\end{itemize}
\end{defn}

\paragraph{Lean.}
\begin{lstlisting}[style=lean]
def directPath : Path G' melbournePort thaiPort :=
  .cons (G := G') sea (.nil (G := G') thaiPort)

def transshipPath : Path G' melbournePort thaiPort :=
  .cons (G := G') sing1 (.cons (G := G') sing2 (.nil (G := G') thaiPort))
\end{lstlisting}

\paragraph{Haskell.}
\begin{lstlisting}[style=haskell]
-- Paths represented as lists of legs (types approximate only)
directPath :: [Leg]
directPath = [Sea]

transshipPath :: [Leg]
transshipPath = [Sing1, Sing2]
\end{lstlisting}


% --- DAG rung theorem -------------------------------------------
\subsection{Theorem: DAG rung (\S7)}

\begin{thm}[DAG rung --- not thin]
The type of paths from \lean{melbournePort} to \lean{thaiPort} in $G'$ is not
a subsingleton: $\lnot\,\mathrm{Subsingleton}(\mathrm{Path}_{G'}(\text{Melbourne
Port},\,\text{Thai Port}))$.

Proof witness: \lean{directPath} has length~1 while \lean{transshipPath} has
length~2, so they are distinct.  Hence the network $\mathcal{N}(G')$ is not
thin at Melbourne Port.
\end{thm}

\paragraph{Lean.}
\begin{lstlisting}[style=lean]
example : Not (Subsingleton (Path G' melbournePort thaiPort)) :=
  fun h => ...  -- length argument; proof omitted
\end{lstlisting}

\paragraph{Haskell.}
\begin{lstlisting}[style=haskell]
-- The type of paths is not a singleton:
-- length directPath /= length transshipPath
dagRung :: Bool
dagRung = length directPath /= length transshipPath  -- True
\end{lstlisting}


%===================================================================
\section{Graph rung: 2-node scrap cycle (\S8)}
%===================================================================

% --- Node / Edge / G_scrap ----------------------------------------
\subsection{Types: \lean{ScrapLoop.Node}, \lean{ScrapLoop.Edge}, \lean{G\_scrap}}

\begin{defn}[Scrap-loop graph]
The scrap-loop graph $G_{\mathrm{scrap}}$ has:
\begin{itemize}
  \item two nodes: \lean{rollingMill} and \lean{meltShop};
  \item two directed edges: \lean{rm2ms} (rolling mill $\to$ melt shop) and
    \lean{ms2rm} (melt shop $\to$ rolling mill).
\end{itemize}
This is a 2-cycle: billets flow from the rolling mill, and scrap returns from
the melt shop.
\end{defn}

\paragraph{Lean.}
\begin{lstlisting}[style=lean]
inductive Node | rollingMill | meltShop
  deriving DecidableEq, Repr

inductive Edge | rm2ms | ms2rm
  deriving DecidableEq, Repr

@[reducible] def G_scrap : Digraph where
  V := Node
  E := Edge
  src := edge_src
  tgt := edge_tgt
\end{lstlisting}

\paragraph{Haskell.}
\begin{lstlisting}[style=haskell]
data Node = RollingMill | MeltShop deriving (Eq, Show)
data Edge = Rm2ms | Ms2rm          deriving (Eq, Show)

gScrap :: Digraph Node Edge
gScrap = Digraph edgeSrc edgeTgt
\end{lstlisting}


% --- oneLap / laps ------------------------------------------------
\subsection{Definitions: \lean{oneLap} and \lean{laps}}

\begin{defn}[Lap paths]
A single lap around the 2-cycle:
\[
  \text{oneLap} : \mathrm{Path}_{G_{\mathrm{scrap}}}(\text{rollingMill},
  \text{rollingMill}),
\]
consisting of the two edges \lean{rm2ms} then \lean{ms2rm}.  The family
\lean{laps} $: \mathbb{N} \to
\mathrm{Path}_{G_{\mathrm{scrap}}}(\text{rollingMill},\text{rollingMill})$
iterates this lap: $\text{laps}(n)$ is the path that traverses the cycle exactly
$n$ times.
\end{defn}

\paragraph{Lean.}
\begin{lstlisting}[style=lean]
def oneLap : Path G_scrap rollingMill rollingMill :=
  .cons (G := G_scrap) rm2ms
    (.cons (G := G_scrap) ms2rm (.nil (G := G_scrap) rollingMill))

def laps : Nat -> Path G_scrap rollingMill rollingMill
  | .zero   => .nil (G := G_scrap) rollingMill
  | .succ n => (laps n).append oneLap
\end{lstlisting}

\paragraph{Haskell.}
\begin{lstlisting}[style=haskell]
oneLap :: [Edge]
oneLap = [Rm2ms, Ms2rm]

laps :: Int -> [Edge]
laps 0 = []
laps n = laps (n-1) ++ oneLap
\end{lstlisting}


% --- Infinite -------------------------------------------------------
\subsection{Structure: \lean{Infinite}}

\begin{defn}[Infinite type (bundled)]
A type $\alpha$ is \textbf{infinite} if there exists an injective function
$\mathbb{N} \to \alpha$:
\[
  \mathrm{Infinite}(\alpha) \;\equiv\; \exists f : \mathbb{N} \to \alpha,\;
  \forall m\, n,\; f(m) = f(n) \Rightarrow m = n.
\]
\end{defn}

\paragraph{Lean.}
\begin{lstlisting}[style=lean]
structure Infinite (alpha : Type) : Prop where
  witness : forall (f : Nat -> alpha), forall m n, f m = f n -> m = n
\end{lstlisting}

\begin{fidelitynote}
The Lean definition quantifies existentially (the witness field records the
injection); the English and Haskell readings make this explicit.
\end{fidelitynote}

\paragraph{Haskell.}
\begin{lstlisting}[style=haskell]
-- A type a is infinite if there is an injection Nat -> a.
-- In Haskell this is the class Infinite from Data.Infinite, or simply:
isInfinite' :: (Int -> a) -> (a -> a -> Bool) -> (Int -> Int -> Bool) -> Bool
isInfinite' f eq implies = True -- just a witness, not a proof
\end{lstlisting}


% --- infinite_pos_rollingMill ------------------------------------
\subsection{Theorem: \lean{infinite\_pos\_rollingMill}}

\begin{thm}[Graph rung --- infinite position set]
The $\Sigma$-type $\sum_{w : \mathrm{Node}} \mathrm{Path}_{G_{\mathrm{scrap}}}(\mathrm{rollingMill},\, w)$
is infinite.  The injection is $n \mapsto (\mathrm{rollingMill},\, \mathrm{laps}(n))$,
which is injective because $\mathrm{laps}(n)$ has length $2n$, a strictly
increasing function of $n$.

Hence the position set of the directed container $\mathcal{N}(G_{\mathrm{scrap}})$
at \lean{rollingMill} is infinite, placing $G_{\mathrm{scrap}}$ on the top rung
of the ladder.
\end{thm}

\paragraph{Lean.}
\begin{lstlisting}[style=lean]
theorem infinite_pos_rollingMill :
    Infinite ((w : Node) * Path G_scrap rollingMill w) where
  witness := forall (n => <|rollingMill, laps n|>), ...
\end{lstlisting}

\paragraph{Haskell.}
\begin{lstlisting}[style=haskell]
-- The injection: n |-> (RollingMill, laps n)
-- Injectivity: length (laps m) = 2*m, so m /= n => laps m /= laps n
injection :: Int -> (Node, [Edge])
injection n = (RollingMill, laps n)
\end{lstlisting}


% --- nontrivialEndo / nontrivialEndo_ne_root ----------------------
\subsection{Definition and theorem: \lean{nontrivialEndo} and \lean{nontrivialEndo\_ne\_root}}

\begin{defn}[Nontrivial endomorphism]
The \textbf{nontrivial endomorphism} at \lean{rollingMill} is the one-lap
position: the pair $(\mathrm{rollingMill},\, \mathrm{oneLap})$ as an element of
$\mathcal{N}(G_{\mathrm{scrap}}).\mathrm{Pos}(\mathrm{rollingMill})$.
\end{defn}

\begin{thm}[Endomorphism is non-trivial]
The nontrivial endomorphism is distinct from the root (the zero-lap empty path):
\[
  \mathrm{nontrivialEndo} \;\neq\;
  \mathcal{N}(G_{\mathrm{scrap}}).\mathrm{root}(\mathrm{rollingMill}).
\]
\end{thm}

\paragraph{Lean.}
\begin{lstlisting}[style=lean]
def nontrivialEndo : (Digraph.network G_scrap).Pos rollingMill :=
  <|rollingMill, oneLap|>

theorem nontrivialEndo_ne_root :
    nontrivialEndo != (Digraph.network G_scrap).root rollingMill
\end{lstlisting}

\paragraph{Haskell.}
\begin{lstlisting}[style=haskell]
nontrivialEndo :: (Node, [Edge])
nontrivialEndo = (RollingMill, oneLap)  -- length 2

root :: (Node, [Edge])
root = (RollingMill, [])  -- length 0

-- nontrivialEndo /= root  (different lengths)
\end{lstlisting}


%===================================================================
\section{Translation conventions}
%===================================================================

\begin{enumerate}
\item \textbf{Paths as lists.}  In Lean, \lean{Path G v w} is an indexed
  inductive type whose constructors carry type-level evidence of endpoint
  compatibility.  In Haskell, paths are approximated as lists of edges
  \hask{[Edge]}; the endpoint type indices are lost.

\item \textbf{Subsingleton.}  Lean's \lean{Subsingleton} is a typeclass
  saying a type has at most one element.  The negation used in the DAG rung
  says: there exist two provably distinct elements.

\item \textbf{Injectivity witness.}  Lean's bundled \lean{Infinite} packs an
  injection together with its injectivity proof.  The Haskell rendering only
  shows the injection function; the proof obligation is dropped.

\item \textbf{Proof bodies omitted.}  Only definitions and theorem
  \emph{statements} are translated.  Lean checks the proofs; the human
  verifier checks that the \emph{statements} match the intended mathematics.
\end{enumerate}

\end{document}
